% AUTOGENERATED DRAFT from manuscript source: manuscript/chapters/00-preface.md
% Converter: scripts/convert_markdown_chapter_to_tex.py
% Review gate: docs/latex-migration-plan.md (Phase 2 per-chapter gate)

\chapter*{写作说明}
\addcontentsline{toc}{chapter}{写作说明}

\section{这本书为谁而写}

如果你会运行一个 PIC case，却还不能回答“粒子在什么时间层被推进”“电流为什么要沿轨迹沉积”“一个小的 \cpath{divE-rho} 残差究竟说明什么”，这本书就是从这里开始的。它不要求读者先熟悉 WarpX 的类层次，也不把源码中的函数名当作知识本身；每个实现细节都先放回物理量、离散方程和可观察结果中解释。

读完主线后，读者应当能够：

\begin{itemize}
\item 从连续方程推导出 PIC 中需要保存和交换的离散量；
\item 沿着一个输入参数追踪到初始化、时间推进、沉积、场更新和诊断输出；
\item 为一个新 case 选择时间步、网格、shape、solver 和诊断，而不是只复制输入文件；
\item 看见 regression 通过时，准确说出它证明了什么、还没有证明什么。
\end{itemize}

\section{如何使用本书}

正文按“模型 -> 离散算法 -> 源码调用链 -> 案例诊断”的顺序展开。章节末的练习要求读者自己定位源码、检查时间层或复现实验；它们不是附加的项目验收清单，而是把阅读变成判断能力的训练。

版本号、运行合同和缺口登记记录的是书稿如何被维护，不是读者必须按时间顺序阅读的内容。正文只在需要解释证据范围时使用它们，不把它们作为学习前提。

本书不是 WarpX 官方文档的翻译，也不是只讲公式的 PIC 理论笔记。它的主线是：先从 Vlasov-Maxwell / Vlasov-Poisson 这类连续模型出发，说明为什么需要宏粒子；再把宏粒子、网格、形函数、沉积和场求解拼成 PIC 算法；最后回到 WarpX 的 \cpath{Source/}、\cpath{Docs/}、\cpath{Examples/} 和 regression 入口，解释一个现代高性能 PIC 程序如何把这些步骤组织成可运行、可扩展、可验证的模拟软件。

源码定位应优先依赖文件职责、函数名和调用关系，而不是固定 commit 或行号。例如，主循环从 \cpath{Evolve()} 与 \cpath{OneStep()} 开始，初始化从 \cpath{InitData()} 开始，粒子/源项主链从 \cpath{PushParticlesandDeposit()} 与 \cpath{SyncCurrentAndRho()} 开始。版本信息只用于复现实验和维护记录，不应改变读者对算法关系的理解。

本书的每个技术判断都应尽量落到六类证据：物理方程、离散公式、WarpX 源码路径、输入参数、示例或测试、文献。DeepWiki、Zread 等 AI 解读页面可以用来快速找到模块名，但不能作为最终依据。

遇到一个新的输入或源码分支时，可按四步阅读：先写出要描述的连续物理量；再标出它在离散网格和时间层上的表示；随后定位 producer 和 consumer；最后选择一个能区分正确与错误的 observable，并写下该观察量不能支持的外推。这样，代码阅读始终服务于物理判断，而不是变成函数名清单。

\subsection{四个跨章节术语}

后续章节反复使用下列词，但它们不是函数名、文件名或测试标签的同义词：

% <-- table block (source: 00-preface L32)
\begin{tabularx}{\textwidth}{LLLL}
\toprule
术语 & 本书中的含义 & 阅读时必须继续追问 & 不能据此直接推出 \\
\midrule
\endfirsthead
\toprule
术语 & 本书中的含义 & 阅读时必须继续追问 & 不能据此直接推出 \\
\midrule
\endhead
source & 被某个场或方程更新实际消费的离散源项或状态，例如同步后的 $\rho$、$\mathbf{J}$ 或谱空间历史量 & 它在哪个时间层、几何位置、AMR level 和同步阶段成立 & 任意名为 \cpath{rho}、\cpath{current\_*} 的数组都已可供 solver 消费 \\
producer & 在确定的生命周期阶段创建、推进、沉积、同步或归约某个状态的代码路径 & 它生成什么量、何时写入、是否还需要过滤、通信或边界处理 & 找到一次写入或对象构造就证明物理状态已经完整 \\
consumer & 读取该状态以更新场、写诊断、作比较或产生后续状态的代码路径 & 它读取的是 producer 的哪一时间层和哪一种表示 & 输出文件成功写出就证明 producer 的物理语义正确 \\
observable & 为回答一个物理问题而定义、可与解析解、守恒关系、reference 或实验量比较的量 & 它的比较对象、容差、几何和不可外推范围是什么 & 任意已输出字段、checksum 或图像都是物理验证 \\
\bottomrule
\end{tabularx}

“证据层级”则描述上述链条有多强：源码可以定位职责，指定案例可以检验给定条件下的 observable，全文文献可以支撑公式或机制；它们必须互相独立，不能把其中任一层放大成其余各层。第 3 章把这四个术语接到控制流和状态交接，第 5--7 章说明 source 怎样经过推进、同步和边界，第 8 章说明 consumer 如何形成诊断，第 9 章再判断外部 reference 的证据强度。

本书默认使用以下记号：粒子位置为 $\mathbf{x}_p$，粒子动量为 $\mathbf{u}_p=\gamma\mathbf{v}_p$，电磁场为 $\mathbf{E},\mathbf{B}$，电荷和电流密度为 $\rho,\mathbf{J}$，粒子权重为 $w_p$，形函数为 $S$。网格量的上标表示时间层，例如 $\mathbf{B}^{n+1/2}$；粒子量一般按 leapfrog 交错在位置和动量时间层上。

\section{按目标选择阅读路线}

不必把全书当作一条只能从头到尾走完的线。先选择你要获得的能力，再沿相应路线完成每个终点练习；三条路线共享术语和证据边界，但不能用其中一条的通过替代另一条的结论。

\subsection{路线 A：先能运行并解释一个最小 PIC case}

适合已经会启动程序、但还不能说明输出为何可信的读者。先读第 1.10--1.14，写出 Debye 长度、时间尺度和一个要保留的 observable；再读第 2.6--2.10，把它放进一个完整 PIC step 和最小输入；接着用第 3.12--3.13 追踪输入、\cpath{warpx\_used\_inputs}、输出与 analysis 的交接；最后完成第 8.15 的 Langmuir 或 uniform-plasma 练习。终点不是“作业跑完”，而是能写出该 case 的 producer、consumer、比较量、容差以及它不能覆盖的 geometry、算法和物理范围。

\subsection{路线 B：先读懂一条从粒子到场的源码链}

适合需要修改或审查实现、但不希望先被所有模块细节淹没的读者。以第 2.7 的导航为索引，读第 3.4--3.10 确定初始化、外层时间步和 source 同步的职责；然后依次完成第 4.16、5.16、6.12 的练习，分别追踪粒子状态、\cpath{rho/J} 和场更新的时间层。每次定位函数时先写“它消费什么、产生什么、下一位 consumer 是谁”，再阅读分派或 kernel；不能从找到一次数组写入就推断求解器已经消费了物理上完整的 source。

\subsection{路线 C：先学会设计验证，而不是只收集输出}

适合准备建立新 case、改变数值算法或解释异常结果的读者。先读第 1.5.1 与第 1.9.1，区分模型选择、能量账本和统计量；再读第 5.14.2.3 与第 5.14.7，把 source residual、解析场、checksum 和收敛趋势拆成不同 consumer；随后以第 8.14--8.16 设计一张诊断记录卡；最后用第 9.8--9.10 复核文献、源码和运行证据能否支持措辞。终点是一份可审查的验证合同，而不是一张看起来合理的场图或一次 regression PASS。

第一次阅读仍可先走路线 A，再按需要进入 B 或 C。无论选择哪条路线，遇到一个新的输入或源码分支，都回到本节前的四步：连续物理量、离散表示、producer/consumer 和带范围说明的 observable。
